For a visual binary stellar system consisting of the stars $\sigma_1$ and
$\sigma_2$, the following relation represents the mass function of the
system:
\[ f(M_1; M_2) = \frac{M_2^3 \sin^3 i}{(M_1 + M_2)^2}, \]
where $M_1$ is the mass of star $\sigma_1$, $M_2$ is the mass of star
$\sigma_2$ and $i$ is the angle between the plane of the stars' orbits and a
plane perpendicular on the direction of observation.

The recorded spectrum of radiations emitted by the star $\sigma_1$, during
several months, reveals a sinusoidal variation of radiation wavelength, with
the period $T = 7$ days and a shift factor
$z = (\Delta\lambda)/\lambda = 0.001$.

\begin{description}
\item[(a)] Prove that the mass function of the system is:
  \[ f(M_1; M_2) = \frac{M_2^3 \sin^3 i}{(M_1 + M_2)^2}
     = \frac{T}{2\pi K} \left( v_1 \cdot \sin i \right)^3, \]
  where $v_1 \cdot \sin i$ is the maximum speed of star $\sigma_1$ relatively
  to the observer; $K$ is the gravitational constant; $i$ is the angle
  between the plane of the orbits and the plane normal to the observation
  direction. Assumptions: the orbits of the stars are circular.
\item[(b)] Derive an expression for the mass function of the system. The
  following values are known: $c = 3 \times 10^8$\,m/s;
  $K = 6.67 \times 10^{-11}\,\mathrm{N\,m^2\,kg^{-2}}$.
\end{description}
